\documentclass[11pt,a4paper]{article}
\usepackage[utf8]{inputenc}
\usepackage[margin=0.9in]{geometry}
\usepackage{graphicx}
\usepackage{amsmath}
\usepackage{amsfonts}
\usepackage{booktabs}
\usepackage{hyperref}
\usepackage{fancyhdr}
\usepackage{titlesec}
\usepackage{float}
\usepackage{caption}
\usepackage{xcolor}

\definecolor{navy}{RGB}{26, 54, 93}
\definecolor{darkgray}{RGB}{45, 55, 72}

\hypersetup{
    colorlinks=true,
    linkcolor=navy,
    filecolor=magenta,      
    urlcolor=cyan,
    pdftitle={Heavy-Lift Multirotor UAV Design Report},
    pdfpagemode=UseOutlines,
}

\pagestyle{fancy}
\fancyhf{}
\rhead{\color{darkgray}\small Heavy-Lift UAV Design Report}
\lhead{\color{darkgray}\small Flight Systems Division}
\rfoot{\color{darkgray}\small Page \thepage}
\lfoot{\color{darkgray}\small Locked Baseline (DESIGN\_LOCK.md)}

\titleformat{\section}{\large\bfseries\color{navy}}{\thesection}{1em}{}
\titleformat{\subsection}{\normalsize\bfseries\color{navy}}{\thesubsection}{1em}{}

\begin{document}

% ==========================================
% COVER PAGE
% ==========================================
\begin{titlepage}
    \centering
    \vspace*{1cm}
    {\Huge\bfseries Design, Aerodynamic BEM, Structural FEA, and Mission Validation of a Heavy-Lift Hybrid Hexacopter\par}
    \vspace{0.8cm}
    {\large\bfseries Comprehensive Engineering Report for a 10 kg Payload, 30 km Operational Mission Range UAV\par}
    \vspace{1.5cm}
    \includegraphics[width=0.75\textwidth]{figures/uav_cad_annotated.png}\par
    \vspace{1.5cm}
    {\large\bfseries Prepared by:\par}
    {\large Flight Systems Engineering Team \par}
    {IIT Madras --- BS Electronic Systems\par}
    \vspace{0.8cm}
    {\large\bfseries Document Control:\par}
    {Locked Design Baseline Reference (\texttt{DESIGN\_LOCK.md})\par}
    \vfill
    {\large Date: July 2026\par}
\end{titlepage}

\newpage
\tableofcontents
\newpage

% ==========================================
% SECTION 1: EXECUTIVE SUMMARY
% ==========================================
\section{Executive Summary \& Mission Requirements}

This technical report documents the complete engineering design, aerodynamic simulation, structural FEA, and electrical integration of a heavy-lift multirotor Unmanned Aerial Vehicle (UAV). The system is engineered to carry a \textbf{10.0 kg payload} over a \textbf{30 km out-and-back operational mission profile} with a \textbf{20-minute on-station hover loiter} and a \textbf{20\% fuel reserve margin}.

To satisfy these aggressive endurance requirements without exceeding physical mass limits, the aircraft utilizes a \textbf{Gas-Electric Hybrid Power Plant} (3.6 kW continuous output driving a 48V DC bus) coupled with a \textbf{Single-Motor Hexacopter Topology} (6 arms at $60^\circ$ radial spacing, 1 motor + 1 ESC + 1 40-inch propeller per arm).

\begin{table}[H]
    \centering
    \caption{Locked System Baseline Specifications (\texttt{DESIGN\_LOCK.md})}
    \begin{tabular}{lll}
        \toprule
        \textbf{Parameter} & \textbf{Locked Value} & \textbf{Engineering Basis / Status} \\
        \midrule
        Aircraft Configuration & Standard Hexacopter & 6 arms @ $60^\circ$ spacing ($0^\circ, 60^\circ, 120^\circ, 180^\circ, 240^\circ, 300^\circ$) \\
        Arm Tube Length ($L$) & \textbf{1.120 m} (1120 mm) & 30mm OD $\times$ 2mm wall carbon tube ($104\text{ mm}$ prop tip clearance) \\
        Propeller Specification & \textbf{40" $\times$ 13"} Carbon Fiber & $R = 0.508\text{ m}$, 2-blade tapered profile, $8.1\text{ g/W}$ hover efficiency \\
        Total Takeoff Weight (TOW) & \textbf{34.575 kg} & Converged via 5-iteration mass-power-fuel loop \\
        Hover Electrical Power & \textbf{3,460.1 W} & Total DC electrical draw ($576.7\text{ W}$ per motor branch @ 48V DC) \\
        Carried Fuel Mass & \textbf{2.323 kg} & 1.859 kg mission fuel burn + 0.464 kg (20\% reserve margin) \\
        Governing Load Case & Asymmetric Motor-Out & 1.5G Emergency Recovery ($\sigma_{\max} = 164.4\text{ MPa}$, $\text{SF} = 4.87 \ge 1.5$) \\
        \bottomrule
    \end{tabular}
\end{table}

% ==========================================
% SECTION 2: CONFIGURATION & GEOMETRY
% ==========================================
\section{Aircraft Configuration \& Geometric Layout}

A rigorous trade study was performed to select the vehicle topology among Quadcopter, Flat Hexacopter, Coaxial Octocopter (X8), and Flat Octocopter configurations:

\begin{enumerate}
    \item \textbf{Single-Motor Hexacopter Selection:} The hexacopter configuration offers single-motor failure recovery capability while eliminating the aerodynamic slipstream interference losses inherent to coaxial pairs (which increase lower rotor hover power by 15--20\%).
    \item \textbf{Arm Length Sizing ($L = 1.12\text{ m}$):} For a 6-arm hexacopter with $60^\circ$ radial spacing, the chord distance between adjacent arm tips equals the arm length ($d = L$). With 40-inch propellers ($D_{\text{prop}} = 1.016\text{ m}$), setting $L = 1.120\text{ m}$ guarantees a tip-to-tip clearance of:
    \begin{equation}
        d_{\text{clearance}} = L - D_{\text{prop}} = 1.120\text{ m} - 1.016\text{ m} = \mathbf{0.104\text{ m}} \quad (104\text{ mm}, 10.2\%\text{ clearance margin})
    \end{equation}
    This eliminates blade tip vortex interference and ensures \textbf{ZERO PROPELLER OVERLAP}.
    \item \textbf{Center Frame Diameter ($400\text{ mm}$):} Twin $3\text{ mm}$ carbon fiber hex plates ($200\text{ mm}$ radius) provide structural clamping for the 6 carbon arms, generator box, fuel tank, and 4 landing gear struts.
\end{enumerate}

\begin{figure}[H]
    \centering
    \includegraphics[width=0.48\textwidth]{figures/uav_cad_topview.png}
    \includegraphics[width=0.48\textwidth]{figures/uav_cad_sideview.png}
    \caption{Top-view rendering displaying 104mm blade tip clearance (left) and Side-view layout displaying 450mm landing gear height and 220mm payload ground clearance (right).}
\end{figure}

% ==========================================
% SECTION 3: AERODYNAMIC BEM & PROPULSION
% ==========================================
\section{Rotor Aerodynamics (BEM Theory) \& Propulsion}

Blade Element Momentum (BEM) theory was implemented in \texttt{rotor\_bem.py} to evaluate the aerodynamic performance of the 40" $\times$ 13" carbon fiber propeller using NACA 4412 airfoil section data across 30 spanwise blade elements.

\subsection{Governing BEM Equations}
The local induced inflow velocity $v_i(r)$ at radial station $r$ is solved numerically via momentum conservation and blade element lift/drag integration:
\begin{equation}
    dT = 4 \pi r \rho (V_\infty + v_i) v_i F dr = \frac{1}{2} N_b \rho W^2 (C_l \cos\phi - C_d \sin\phi) c(r) dr
\end{equation}
where $N_b = 2$ blades, $\phi = \arctan\left(\frac{V_\infty + v_i}{\Omega r}\right)$ is the local inflow angle, $W = \sqrt{(\Omega r)^2 + (V_\infty + v_i)^2}$ is the resultant flow velocity, and $F(r)$ is Prandtl's tip loss factor:
\begin{equation}
    F(r) = \frac{2}{\pi} \arccos\left( \exp\left[ -\frac{N_b (R - r)}{2 r \sin\phi} \right] \right)
\end{equation}

\subsection{Hover Efficiency Verification}
At the target Takeoff Weight of $34.575\text{ kg}$, each of the 6 rotors must produce a hover thrust of:
\begin{equation}
    T_{\text{rotor}} = \frac{34.575\text{ kg} \times 9.81\text{ m/s}^2}{6} = \mathbf{56.53\text{ N}} \quad (5.76\text{ kg})
\end{equation}
The BEM simulation confirms:
\begin{itemize}
    \item \textbf{Target Hover Speed:} $1,850\text{ RPM}$
    \item \textbf{Figure of Merit (FoM):} $0.72$
    \item \textbf{Rotor Hover Efficiency:} $8.1\text{ g/W}$ ($547.2\text{ W}$ mechanical power per rotor)
    \item \textbf{Total Electrical Power Draw:} $3,460.1\text{ W}$ (including $94.6\%$ motor/ESC efficiency and $150\text{ W}$ auxiliary avionics).
\end{itemize}

% ==========================================
% SECTION 4: MASS BUDGET & INERTIA TENSOR
% ==========================================
\section{Converged Mass Budget, CG, and Inertia Tensor}

A 5-iteration dynamic mass-power-fuel convergence loop (\texttt{run\_tow\_convergence}) was executed in \texttt{mass\_budget.py}. The system converged to a Total Takeoff Weight of \textbf{34.575 kg}.

\begin{table}[H]
    \centering
    \caption{Converged System Mass Budget Breakdown}
    \begin{tabular}{llccc}
        \toprule
        \textbf{Subsystem Category} & \textbf{Component Description} & \textbf{Qty} & \textbf{Unit Mass (kg)} & \textbf{Total Mass (kg)} \\
        \midrule
        Payload Subsystem & Cargo Bay Box / Gimbal Placeholder & 1 & 10.000 & 10.000 \\
        Power Plant & Hybrid Generator (Dry 3.6 kW Engine) & 1 & 4.500 & 4.500 \\
        Power Plant & Gasoline Fuel (Carried with 20\% Reserve) & 1 & 2.323 & 2.323 \\
        Power Plant & 12S LiPo Buffer Battery Pack & 1 & 1.500 & 1.500 \\
        Propulsion System & T-Motor U15 II KV100 Brushless Motors & 6 & 1.050 & 6.300 \\
        Propulsion System & 40" $\times$ 13" Carbon Fiber Propellers & 6 & 0.250 & 1.500 \\
        Propulsion System & T-Motor Flame 120A FOC ESCs & 6 & 0.150 & 0.900 \\
        Structural Frame & Center Carbon Hex Plates (400mm Dia) & 2 & 1.100 & 2.200 \\
        Structural Frame & Carbon Arm Tubes (30mm OD $\times$ 1.12m) & 6 & 0.392 & 2.352 \\
        Structural Frame & Landing Gear Struts \& Skid Pads & 4 & 0.300 & 1.200 \\
        Avionics \& Signal & Flight Controller, GPS/RTK, Antenna, Wiring & 1 & 1.800 & 1.800 \\
        \midrule
        \textbf{TOTAL TAKEOFF WEIGHT (TOW)} & & & & \textbf{34.575 kg} \\
        \bottomrule
    \end{tabular}
\end{table}

\subsection{Center of Gravity \& Inertia Tensor Calculation}
Relative to the frame geometric origin $[0, 0, 0]$, the Center of Gravity (CG) is located at:
\begin{equation}
    \mathbf{r}_{\text{CG}} = [-0.0000, -0.0006, -0.0533]\text{ m}
\end{equation}
The mass moment of inertia tensor $\mathbf{I}_{\text{CG}}$ about the CG is calculated via parallel axis theorem $\mathbf{I} = \sum m_i \left( (\mathbf{r}_i \cdot \mathbf{r}_i) \mathbf{E} - \mathbf{r}_i \otimes \mathbf{r}_i \right)$:
\begin{equation}
    \mathbf{I}_{\text{CG}} = \begin{bmatrix}
        6.187 & 0.000 & 0.000 \\
        0.000 & 6.168 & 0.000 \\
        0.000 & 0.000 & 11.670
    \end{bmatrix} \text{ kg}\cdot\text{m}^2
\end{equation}

% ==========================================
% SECTION 5: MISSION ENERGY & ENDURANCE
% ==========================================
\section{Operational Mission Profile \& Energy Simulation}

The vehicle operational profile was simulated in \texttt{power\_endurance.py} over a 30 km out-and-back mission profile:
\begin{enumerate}
    \item \textbf{Climb Phase:} 100 m vertical climb @ 2 m/s ($50\text{ s}$, $4,200\text{ W}$ power).
    \item \textbf{Cruise Out Phase:} 30 km forward cruise @ 12 m/s ($2,500\text{ s} = 41.7\text{ min}$, $2,450\text{ W}$ power with translational lift benefit).
    \item \textbf{On-Station Loiter Phase:} 20 min stationary hover ($1,200\text{ s}$, $3,460\text{ W}$ power).
    \item \textbf{Cruise Back Phase:} 30 km return cruise @ 12 m/s ($2,500\text{ s} = 41.7\text{ min}$, $2,450\text{ W}$ power).
\end{enumerate}

\begin{figure}[H]
    \centering
    \includegraphics[width=0.75\textwidth]{figures/endurance_simulation.png}
    \caption{Dynamic mission power draw and fuel mass burn profile over the 30 km out-and-back mission.}
\end{figure}

The total mission fuel burn is \textbf{1.859 kg}. Applying the required $20\%$ fuel reserve margin ($\text{Fuel}_{\text{carried}} = \text{Fuel}_{\text{burn}} / 0.8$) dictates a carried fuel mass of \textbf{2.323 kg}, providing \textbf{20.8 minutes of reserve hover time}.

% ==========================================
% SECTION 6: STRUCTURAL ANALYSIS & FEA
% ==========================================
\section{Structural FEA Analysis \& Dual Load Cases}

The 1.12 m carbon fiber arm tubes (30 mm OD, 2 mm wall thickness, $I_{\text{area}} = 1.706 \times 10^{-8}\text{ m}^4$, $\text{UTS} = 800\text{ MPa}$, $E = 120\text{ GPa}$) were evaluated under two independent load cases in \texttt{structural\_analysis.py}.

\begin{table}[H]
    \centering
    \caption{Structural Load Case Evaluation \& Governing Case Selection}
    \begin{tabular}{lcc}
        \toprule
        \textbf{Structural Evaluation Parameter} & \textbf{Case 1: Symmetric 2.5G Limit} & \textbf{Case 2: Asymmetric Motor-Out (Governing)} \\
        \midrule
        Total Vertical Lift Force & $847.9\text{ N}$ & $508.8\text{ N}$ \\
        Effective Active Arm Count & 6 arms & 3 arms (moment balanced under recovery) \\
        Peak Force per Active Arm & $141.3\text{ N}$ & \textbf{169.6 N (+20.0\% higher load)} \\
        Root Bending Moment ($M = F \times L$) & $158.3\text{ N-m}$ & \textbf{189.9 N-m} \\
        Peak Bending Stress ($\sigma_{\max} = \frac{M c}{I}$) & $137.0\text{ MPa}$ & \textbf{164.4 MPa} \\
        Tip Deflection ($\delta = \frac{F L^3}{3 E I}$) & $11.60\text{ mm}$ & \textbf{13.92 mm} \\
        Safety Factor ($\text{SF} = \text{UTS} / \sigma_{\max}$) & $\text{SF} = 5.84$ & \textbf{SF = 4.87 (PASSED $\ge 1.5$)} \\
        \bottomrule
    \end{tabular}
\end{table}

\textbf{Governing Load Case Selection:} The \textbf{Asymmetric Motor-Out Emergency Recovery Case GOVERNS} the structural sizing because motor failure compensation concentrates higher force ($169.6\text{ N}$ vs $141.3\text{ N}$) on active arms adjacent to the failed rotor, producing the peak bending stress of \textbf{164.4 MPa} ($\text{SF} = 4.87 > 1.5$).

\begin{figure}[H]
    \centering
    \includegraphics[width=0.75\textwidth]{figures/arm_structural_fea.png}
    \caption{Euler-Bernoulli FEA beam bending stress and vertical deflection along the 1.12m carbon arm.}
\end{figure}

% ==========================================
% SECTION 7: ELECTRICAL & CAD DELIVERABLES
% ==========================================
\section{Electrical Architecture \& 3D CAD Deliverables}

\subsection{Electrical Schematic (KiCad Baseline)}
The electrical distribution system (\texttt{power\_and\_signal\_schematic.kicad\_sch}) is designed around a 48V DC bus supplied by the 3.6 kW hybrid generator. Total hover current draw is $72.1\text{ A}$ ($12.02\text{ A}$ per ESC branch), protected by 120A Flame FOC ESCs and a 100A BMS battery failsafe coupling.

\subsection{Parametric 3D CAD Assembly (FreeCAD \& STEP)}
A high-fidelity parametric CAD assembly macro (\texttt{hexacopter\_assembly.py}) and ISO STEP AP214 file (\texttt{hexacopter\_assembly.step}) were generated containing 34 distinct subsystem components (including hollow arm tubes, two-tier motors, 4x M4 bolt patterns, 2-blade tapered props, $300 \times 200 \times 150\text{ mm}$ cargo bay box, EO camera gimbal with lens barrel, telemetry antenna, and landing gear skid pads).

\begin{figure}[H]
    \centering
    \includegraphics[width=0.9\textwidth]{figures/uav_cad_annotated.png}
    \caption{Annotated 3D Technical CAD Assembly displaying all core multirotor subsystems.}
\end{figure}

% ==========================================
% SECTION 8: REQUIREMENTS TRACEABILITY
% ==========================================
\section{Deliverables Requirements Traceability Matrix}

\begin{table}[H]
    \centering
    \caption{Assignment Deliverables \& Design Lock Compliance Matrix}
    \begin{tabular}{p{3.5cm}p{4cm}p{4.5cm}l}
        \toprule
        \textbf{Deliverable Item} & \textbf{Target Specification} & \textbf{Verification Method / Code} & \textbf{Status / Result} \\
        \midrule
        1. Configuration Fix & 6-arm single-motor hexacopter & \texttt{mass\_budget.py}, \texttt{DESIGN\_LOCK.md} & \textbf{PASSED:} 6 arms @ $60^\circ$ \\
        2. Propeller Sizing & 40" $\times$ 13" ($R = 0.508\text{ m}$) & \texttt{rotor\_bem.py}, \texttt{generate\_cad\_model.py} & \textbf{PASSED:} 104mm prop clearance \\
        3. Mission Energy & 30 km out-and-back + 20m loiter & \texttt{power\_endurance.py} & \textbf{PASSED:} 2.323 kg fuel (20\% res) \\
        4. Load Cases & Independent 2.5G \& motor-out & \texttt{structural\_analysis.py} & \textbf{PASSED:} Motor-out governs (SF=4.87) \\
        5. TOW Convergence & 5-iteration convergence loop & \texttt{run\_tow\_convergence} & \textbf{PASSED:} TOW = 34.575 kg \\
        6. CG \& Inertia & Recalculated for 6-arm layout & \texttt{calculate\_inertia\_tensor} & \textbf{PASSED:} $I_{zz} = 11.670\text{ kg-m}^2$ \\
        7. Deliverables & FreeCAD, STEP, XFLR5, KiCad & \texttt{hexacopter\_assembly.py}, \texttt{.kicad\_sch} & \textbf{PASSED:} 100\% generated \\
        \bottomrule
    \end{tabular}
\end{table}

\end{document}
